\documentclass[11pt,a4paper]{article}

% ---------------------------------------------------------
% PREAMBLE (stable, warning-free, Overleaf-safe)
% ---------------------------------------------------------
\usepackage[utf8]{inputenc}
\usepackage[T1]{fontenc}
\usepackage{lmodern}
\usepackage{amsmath, amssymb, amsfonts}
\usepackage{bm}
\usepackage{geometry}
\usepackage{graphicx}
\usepackage{hyperref}
\usepackage{cite}
\usepackage{authblk}
\usepackage{physics}
\usepackage{mathtools}

\geometry{margin=2.4cm}

\title{\textbf{Companion CI V2.0: Black Holes and White Holes as CPT-Conjugate Inter-Sector Horizons in RDCC}}
\author{Michael Lehmann \\ \texttt{mi.lehmann@gmx.de}}
\date{May 2026}

\begin{document}
\maketitle

\begin{abstract}
In the maximally extended Schwarzschild solution a white hole is the time-reversed counterpart of a black hole. Within the Relaxation-Driven Cyclic Cosmology (RDCC) this mathematical duality acquires a physical interpretation. Because the global cosmological state is CPT-symmetric and resides in a bipartite Hilbert space, the mirror sector $(-)$ is the exact CPT-conjugate of the visible sector $(+)$. A black hole in the $(+)$ sector appears as a white hole in the $(-)$ sector, and vice versa. The apparent one-way causal structure of horizons is therefore a projection effect of sectoral observation rather than an absolute property of spacetime.

The minimal dimension-six inter-sector operator and the relaxation dynamics governed by the parameter $\alpha$ soften classical singularities, regulate horizon microphysics, and ensure global information conservation. Black holes and white holes emerge as two sectoral projections of a single CPT-symmetric inter-sector horizon. This Companion develops the unified RDCC interpretation of horizon physics and outlines observational consequences for gravitational-wave echoes, the stochastic gravitational-wave background, and weak lensing around supermassive black holes.
\end{abstract}

% ---------------------------------------------------------
% 1. INTRODUCTION
% ---------------------------------------------------------

\section{Introduction}

Black holes are among the most robust predictions of General Relativity. Observations of X-ray binaries, gravitational-wave merger events, and horizon-scale imaging have firmly established their astrophysical reality. White holes, which appear in the maximally extended Schwarzschild solution as time-reversed counterparts of black holes, are usually regarded as unphysical artifacts of analytic continuation.

The Relaxation-Driven Cyclic Cosmology (RDCC) framework offers a different perspective. RDCC postulates that the global cosmological state lives in a bipartite Hilbert space and is exactly invariant under a global CPT transformation. Observers are confined to one sector and only have access to a sectoral projection of the global state. Causal structure, entropy flow, and the arrow of time are therefore relative to the chosen sector.

In this setting the black-hole/white-hole duality of the maximally extended Schwarzschild geometry acquires a natural physical interpretation. A black hole in the $(+)$ sector is mapped by CPT to a white hole in the $(-)$ sector. The one-way nature of horizons is not an absolute feature of spacetime but a consequence of sectoral projection. From the global point of view the horizon is an inter-sector interface across which information is redistributed but not destroyed.

This Companion develops this interpretation systematically. We review the Schwarzschild geometry and its CPT structure, embed it into the bipartite Hilbert space of RDCC, analyze the role of the inter-sector operator and relaxation dynamics near horizons, and show how global unitarity resolves the black-hole information paradox. We conclude with observational implications for gravitational-wave echoes, the stochastic gravitational-wave background, and weak lensing around supermassive black holes.
% ---------------------------------------------------------
% 2. RDCC FRAMEWORK AND BIPARTITE HILBERT SPACE
% ---------------------------------------------------------

\section{RDCC framework and bipartite Hilbert space}

RDCC is built on a single structural postulate: the complete cosmological quantum state 
\(\ket{\Psi_{\rm global}}\) lives in a bipartite Hilbert space
\begin{equation}
\mathcal{H}_{\rm global} = \mathcal{H}_{+} \otimes \mathcal{H}_{-},
\end{equation}
and is exactly invariant under a global CPT transformation that exchanges the two sectors. 
The visible sector \((+)\) and the mirror sector \((-)\) are related by
\begin{equation}
\ket{\Psi_{\rm global}} = \widehat{\mathrm{CPT}} \ket{\Psi_{\rm global}}, 
\qquad 
\widehat{\mathrm{CPT}} : \mathcal{H}_{+} \leftrightarrow \mathcal{H}_{-}.
\end{equation}

Observers are confined to one sector and only have access to the reduced density matrix obtained by tracing out the other sector. For an observer in the \((+)\) sector the relevant state is
\begin{equation}
\rho_{+} = \mathrm{Tr}_{-} \left[ \ket{\Psi_{\rm global}} \bra{\Psi_{\rm global}} \right].
\end{equation}
The thermodynamic arrow of time, entropy increase, and the apparent irreversibility of macroscopic processes arise from this sectoral projection combined with relaxation dynamics that drive the reduced state toward equilibrium.

The same structure applies not only to cosmological degrees of freedom but also to local gravitational configurations such as black holes. A black hole in the \((+)\) sector is part of the global state and has a CPT-conjugate counterpart in the \((-)\) sector. The horizon is then naturally interpreted as an inter-sector interface rather than a purely intra-sector boundary.

RDCC introduces a minimal inter-sector coupling through a dimension-six operator of the schematic form
\begin{equation}
\mathcal{O}_{\rm int} \sim \frac{1}{M^{2}} \, T^{(+)}_{\mu\nu} \, T^{(-)\mu\nu},
\end{equation}
where \(T^{(\pm)}_{\mu\nu}\) are the energy-momentum tensors in the two sectors and \(M\) is a high-energy scale. This operator mediates weak but nonzero correlations between the sectors. The strength and effective impact of this coupling are controlled by a relaxation parameter \(\alpha\), which also governs the cyclic cosmological dynamics.

In the cosmological context this structure leads to a nonsingular bounce, a natural resolution of the Tolman low-entropy problem, and correlated small deviations from the standard cosmological model. In the present Companion we apply the same structural ingredients to black-hole spacetimes.

% ---------------------------------------------------------
% 3. SCHWARZSCHILD GEOMETRY AND CPT TIME REVERSAL
% ---------------------------------------------------------

\section{Schwarzschild geometry and CPT time reversal}

The exterior Schwarzschild metric in standard coordinates is
\begin{equation}
ds^{2} = -\left(1 - \frac{2M}{r}\right) dt^{2} 
+ \left(1 - \frac{2M}{r}\right)^{-1} dr^{2} 
+ r^{2} d\Omega^{2},
\end{equation}
where \(M\) is the mass parameter and \(d\Omega^{2}\) is the metric on the unit two-sphere. 
This coordinate patch covers the region outside the event horizon at \(r = 2M\).

To understand the full causal structure one introduces Kruskal–Szekeres coordinates \((U,V)\), 
in which the maximally extended Schwarzschild solution is analytic and contains four regions:
the exterior region I, the black-hole interior II, the white-hole interior III, and a second exterior region IV. 
The event horizons are null surfaces separating these regions.

Time reversal in this context exchanges the black-hole and white-hole regions. 
A black hole is characterized by a future event horizon from which no signal can escape, 
while a white hole has a past event horizon into which no signal can enter. 
In the purely classical theory both structures are allowed by the field equations, 
but only black holes are considered physically realized.

From the point of view of RDCC this asymmetry is not fundamental. 
The global CPT symmetry of the full cosmological state implies that for every black-hole configuration in the \((+)\) sector there exists a CPT-conjugate configuration in the \((-)\) sector. 
In the maximally extended Schwarzschild geometry this conjugate configuration is naturally identified with the white-hole region.

The full Kruskal diagram can therefore be reinterpreted as a bipartite structure: 
one exterior region and its associated black-hole interior belong to the \((+)\) sector, 
while the other exterior region and its associated white-hole interior belong to the \((-)\) sector. 
The two horizons are sectoral projections of a single CPT-symmetric inter-sector interface.
% ---------------------------------------------------------
% 4. SECTORAL PROJECTION AND RELATIVE CAUSALITY
% ---------------------------------------------------------

\section{Sectoral projection and relative causality}

In RDCC observers are confined to one sector and only have access to the sectoral projection 
of the global state. This has profound consequences for the interpretation of causal structure.

Consider a global configuration in which the horizon is an inter-sector interface. 
From the point of view of the global state, information can flow across this interface in a way 
that preserves unitarity and CPT symmetry. However, an observer in the \((+)\) sector only sees 
the projection of this process onto \(\mathcal{H}_{+}\). For such an observer the horizon appears as a 
one-way surface: nothing can escape from the black-hole interior back to the exterior region.

From the perspective of the \((-)\) sector the same global process appears time-reversed. 
The object that corresponds to the black hole in the \((+)\) sector is seen as a white hole in the 
\((-)\) sector. For an observer in the \((-)\) sector nothing can enter this object from the exterior, 
but signals can emerge from it. The causal structure is therefore relative to the sectoral projection.

This leads to the following picture:
\begin{itemize}
    \item In the \((+)\) sector the horizon is a future event horizon of a black hole.
    \item In the \((-)\) sector the same global interface appears as a past event horizon of a white hole.
\end{itemize}

The one-way nature of horizons is not an absolute property of spacetime but a consequence of 
restricting attention to a single sector. The global state sees a symmetric interface, while each 
sector sees an asymmetric horizon.

This interpretation is fully consistent with the RDCC view of the thermodynamic arrow of time. 
Entropy increase and irreversibility are sectoral phenomena that arise from projection and relaxation. 
The same mechanism that produces a time-oriented cosmology in each sector also produces a 
time-oriented horizon structure.

% ---------------------------------------------------------
% 5. INTER-SECTOR OPERATOR AND RELAXATION AT HORIZONS
% ---------------------------------------------------------

\section{Inter-sector operator and relaxation at horizons}

The minimal inter-sector operator
\begin{equation}
\mathcal{O}_{\rm int} \sim \frac{1}{M^{2}} \, 
T^{(+)}_{\mu\nu} \, T^{(-)\mu\nu}
\end{equation}
couples the energy-momentum content of the two sectors. 
Near a black-hole horizon the local energy densities and curvature invariants can become large, 
which enhances the effective impact of this operator. 
The relaxation dynamics governed by the parameter \(\alpha\) then play a crucial role.

We can distinguish three regimes:

\begin{enumerate}
    \item \textbf{Far from the horizon:}  
    The inter-sector coupling is extremely weak, and the spacetime is well described by the 
    classical Schwarzschild solution in each sector. 
    The sectors are effectively decoupled at the level of local observables.

    \item \textbf{Near the horizon:}  
    The combination of strong gravity and nonzero inter-sector coupling leads to enhanced 
    entanglement between the sectors. 
    The relaxation dynamics drive the local state toward a configuration in which the classical 
    singularity is replaced by a smooth, nonsingular region. 
    The horizon becomes an inter-sector interface across which information is redistributed.

    \item \textbf{Deep in the interior:}  
    The classical singularity is avoided. 
    Instead, the global state undergoes a bounce-like transition in which degrees of freedom 
    are reconfigured between the sectors. 
    The relaxation timescale
    \begin{equation}
        \tau_{\rm rel} \sim \frac{1}{\alpha}
    \end{equation}
    controls how long the black hole appears quasi-eternal from the perspective of a sectoral observer.
\end{enumerate}

The key point is that the inter-sector operator does not introduce arbitrary new fields or 
ad-hoc boundary conditions. 
It is the same structural ingredient that appears in the cosmological context. 
The role of \(\alpha\) is also unified: it governs both the large-scale relaxation that drives the 
cyclic cosmology and the local relaxation that smooths out horizon singularities.
% ---------------------------------------------------------
% 6. GLOBAL INFORMATION CONSERVATION AND THE BLACK-HOLE INFORMATION PARADOX
% ---------------------------------------------------------

\section{Global information conservation and the black-hole information paradox}

In standard semiclassical gravity the black-hole information paradox arises because Hawking 
radiation appears thermal and uncorrelated with the initial state. If the black hole evaporates 
completely, the final state seems to be mixed even if the initial state was pure. This conflicts 
with the unitarity of quantum mechanics.

In RDCC the fundamental object is the global state \(\ket{\Psi_{\rm global}}\) in the bipartite 
Hilbert space. This state evolves unitarily under a CPT-symmetric Hamiltonian that includes 
the inter-sector operator. Information is never lost at the global level. What appears as 
information loss in one sector is compensated by information flow into the other sector.

From the perspective of an observer in the \((+)\) sector the formation and evaporation of a 
black hole look similar to the standard picture. Matter collapses, a horizon forms, and 
Hawking-like radiation is emitted. However, the detailed structure of the radiation is modified 
by the inter-sector coupling. Correlations that would be invisible in a purely intra-sector 
description are now encoded in subtle entanglement patterns between the emitted quanta and 
degrees of freedom in the \((-)\) sector.

The apparent loss of information in the \((+)\) sector is therefore an artifact of tracing out the 
\((-)\) sector. If one had access to the full global state, one would see that the process is 
unitary. The black-hole information paradox is thus reinterpreted as a \emph{projection paradox}: 
it arises from insisting on a single-sector description of a fundamentally bipartite process.

This perspective also clarifies the role of white holes. In the global picture the black-hole 
horizon in the \((+)\) sector and the white-hole horizon in the \((-)\) sector are two aspects of 
the same inter-sector interface. Information that appears to fall irreversibly into the black hole 
from the \((+)\) perspective can re-emerge in the \((-)\) sector as white-hole-like emission. 
The global CPT symmetry ensures that the full process is time-symmetric and unitary, even 
though each sector sees a time-oriented, apparently irreversible evolution.

A further structural consequence of the bipartite Hilbert-space formulation is that a black hole 
in the \((+)\) sector and its CPT-conjugate white hole in the \((-)\) sector do not represent two 
independent objects. They form a single global configuration whose apparent separation arises 
only from sectoral projection. From the viewpoint of the global state the horizon is therefore 
not a terminal boundary but an inter-sector interface across which information is redistributed 
in a CPT-symmetric manner.

This implies that the black-hole/white-hole pair acts as a global information bridge between 
the two sectors. Energy and information are conserved globally, but each sector sees only its 
own projection of the process. What appears as irreversible infall in the \((+)\) sector corresponds 
to time-reversed outflow in the \((-)\) sector. The two processes are complementary aspects of a 
single CPT-symmetric evolution.

% ---------------------------------------------------------
% 7. OBSERVATIONAL IMPLICATIONS
% ---------------------------------------------------------

\section{Observational implications}

Although the inter-sector coupling is weak, the combination of strong gravity and long timescales 
near black holes opens the possibility of observable signatures. We outline three classes of 
potential effects.

\subsection{Gravitational-wave echoes}

If the horizon is not a perfectly absorbing surface but an inter-sector interface with nontrivial 
microphysics, the ringdown signal of a perturbed black hole may exhibit late-time echoes. 
In RDCC the relaxation dynamics near the horizon introduce a characteristic timescale 
\(\tau_{\rm rel} \sim 1/\alpha\). This timescale can manifest itself as a delay between the primary 
ringdown and subsequent echo-like features.

The detailed waveform depends on the effective reflectivity of the inter-sector interface and the 
frequency dependence of the inter-sector coupling. Qualitatively, one expects small-amplitude, 
phase-shifted repetitions of the ringdown signal at intervals set by the light-travel time between 
the effective interface and the photon sphere, modulated by the relaxation dynamics. 
Future observations with LIGO, Virgo, KAGRA, LISA, and the Einstein Telescope could 
constrain such scenarios.

\subsection{Stochastic gravitational-wave background}

The cumulative effect of many black-hole formation and merger events in the presence of 
inter-sector coupling can imprint characteristic features on the stochastic gravitational-wave 
background. If horizon microphysics in RDCC modifies the late-time tail of ringdown signals or 
introduces additional channels of energy transfer between sectors, the resulting background may 
deviate from the standard prediction.

In particular, the same relaxation parameter \(\alpha\) that governs cosmological deviations from 
the standard model can induce a mild scale dependence in the gravitational-wave spectrum. 
This connects the black-hole sector to the cosmological sector and provides a cross-check of 
the RDCC framework.

\subsection{Weak lensing around supermassive black holes}

Inter-sector leakage near horizons can in principle alter the effective stress-energy distribution 
in the vicinity of supermassive black holes. While the effect is expected to be small, it could lead 
to subtle deviations in the lensing patterns of background sources. High-resolution imaging and 
precise weak-lensing measurements around galactic centers might reveal small anomalies that 
are difficult to explain within a purely classical framework.

Such signatures would be challenging to isolate, but they illustrate that RDCC does not confine 
its predictions to the early universe. The same structural postulate that shapes the global 
cosmology also has local consequences in strong-gravity regimes.
% ---------------------------------------------------------
% 8. RELATION TO THE THERMODYNAMIC ARROW OF TIME AND CYCLIC COSMOLOGY
% ---------------------------------------------------------

\section{Relation to the thermodynamic arrow of time and cyclic cosmology}

One of the central achievements of RDCC is the unification of the thermodynamic arrow of time, 
the resolution of cosmological singularities, and the cyclic structure of the universe under a 
single structural postulate: global CPT symmetry of the bipartite cosmological quantum state 
combined with relaxation dynamics.

The present Companion shows that the same ingredients naturally extend to black-hole physics. 
The apparent irreversibility of black-hole formation and evaporation is a sectoral phenomenon, 
just like the cosmological arrow of time. The classical singularity inside a black hole is replaced 
by a nonsingular region in which the global state undergoes a bounce-like reconfiguration 
between sectors, analogous to the cosmological bounce.

In this sense black holes are local echoes of the global cyclic structure. They are regions where 
the interplay between CPT symmetry, bipartite structure, and relaxation dynamics becomes 
manifest on astrophysical scales. The black-hole information paradox is then not an isolated 
puzzle but part of a broader pattern: whenever one insists on a single-sector description of a 
fundamentally bipartite, CPT-symmetric process, apparent paradoxes arise that dissolve once 
the global structure is taken into account.

% ---------------------------------------------------------
% 9. CONCLUSION
% ---------------------------------------------------------

\section{Conclusion}

In this Companion we have developed a RDCC-based interpretation of black holes and white holes 
as CPT-conjugate inter-sector horizons. Starting from the maximally extended Schwarzschild 
solution, we have shown that the mathematical duality between black and white holes acquires 
a physical meaning in a bipartite, CPT-symmetric cosmological framework. A black hole in the 
visible \((+)\) sector corresponds to a white hole in the mirror \((-)\) sector, and the one-way 
nature of horizons is a projection effect of sectoral observation.

The minimal dimension-six inter-sector operator and the relaxation dynamics governed by the 
parameter \(\alpha\) play a central role. Near horizons they induce entanglement leakage between 
sectors, smooth out classical singularities, and ensure global information conservation. The 
black-hole information paradox is reinterpreted as a projection paradox: information is never 
lost at the global level, but sectoral observers who trace out the other sector perceive an 
apparent loss.

We have outlined possible observational signatures in gravitational-wave echoes, the stochastic 
gravitational-wave background, and weak lensing around supermassive black holes. While these 
effects are subtle, they illustrate that RDCC connects cosmological and astrophysical phenomena 
within a single structural framework.

This Companion thus extends the reach of RDCC from the largest scales of the universe down 
to the strong-gravity regime of black holes. It reinforces the central message of the program: 
many long-standing puzzles in cosmology and gravity can be addressed coherently by combining 
global CPT symmetry, a bipartite Hilbert space, and relaxation dynamics with a single free 
parameter.

% ---------------------------------------------------------
% REFERENCES
% ---------------------------------------------------------

\bibliographystyle{apsrev4-2}
\nocite{*}
\bibliography{references}

\end{document}
